\section*{How the Proof Works}

The global three-dimensional Navier--Stokes regularity problem remains open in
public mathematics. This section explains the route claimed by this manuscript.

The basic energy identity is
\[
\frac12\|u(t)\|_2^2
+\nu\int_0^t\|\nabla u(\tau)\|_2^2\,d\tau
=\frac12\|u_0\|_2^2.
\]
It keeps the total kinetic energy finite and controls the gradient after
integration in time. It does not bound the gradient at each instant.

A stretching vortex shows why that gap matters. Incompressibility narrows the
core as the tube lengthens. The velocity then changes across less space, so the
gradient can become much larger without forcing the total kinetic energy to
blow up.

This is why scaling enters. It tells us whether a compressed structure can
still have the Navier--Stokes form. From a solution \((u,p)\), form
\[
u_\lambda(x,t)=\lambda u(\lambda x,\lambda^2t),
\qquad
p_\lambda(x,t)=\lambda^2p(\lambda x,\lambda^2t)
\qquad (\lambda>1).
\]
A structure of width \(r\) is squeezed to width \(r/\lambda\), its velocity is
multiplied by \(\lambda\), and its time scale is divided by \(\lambda^2\).
Differentiating the rescaled fields gives
\[
\begin{aligned}
\partial_tu_\lambda(x,t)
&=\lambda^3(\partial_tu)(\lambda x,\lambda^2t),\\
(u_\lambda\!\cdot\nabla)u_\lambda(x,t)
&=\lambda^3((u\!\cdot\nabla)u)(\lambda x,\lambda^2t),\\
\nabla p_\lambda(x,t)
&=\lambda^3(\nabla p)(\lambda x,\lambda^2t),\\
\nu\Delta u_\lambda(x,t)
&=\lambda^3\nu(\Delta u)(\lambda x,\lambda^2t).
\end{aligned}
\]
Every term picks up \(\lambda^3\), so the compressed pair is again a solution.
This compares two solutions at different scales; it does not say that one
vortex automatically evolves through a chain of rescaled copies.

The effect on spatial Sobolev norms follows from
\[
\widehat{u_\lambda}(\xi,t)
=\lambda^{-2}\widehat u(\xi/\lambda,\lambda^2t):
\]
\[
\begin{aligned}
\|u_\lambda(\cdot,t)\|_{\dot H^s}^2
&=\int_{\mathbb R^3}|\xi|^{2s}\lambda^{-4}
  |\widehat u(\xi/\lambda,\lambda^2t)|^2\,d\xi\\
&=\lambda^{2s-1}
  \|u(\cdot,\lambda^2t)\|_{\dot H^s}^2.
\end{aligned}
\]
For \(s<\tfrac12\), the factor falls as \(\lambda\) grows. For
\(s>\tfrac12\), it rises. At \(s=\tfrac12\), it is exactly one: concentration
changes the scale of the flow without changing its \(\dot H^{1/2}\) size. At
\(s=0\), the kinetic energy scales by \(\lambda^{-1}\). It can stay finite
while every norm above the critical height grows.

If the core contracts by a fixed factor \(0<\rho<1\) at each turn, iterating
the rescaling gives
\[
r_j=\rho^jr_0,
\qquad
U_j=\rho^{-j}U_0,
\qquad
\tau_j=\rho^{2j}\tau_0.
\]
Each turn takes \(\rho^2\) times as long as the one before it, so the durations
form a convergent geometric series:
\[
\sum_{j=0}^{\infty}\tau_j
=
\tau_0\sum_{j=0}^{\infty}\rho^{2j}
=
\frac{\tau_0}{1-\rho^2}<\infty.
\]
That is the Zeno danger: the velocity scale rises through infinitely many
thinner cores while the elapsed time stays finite. This does not prove that
Navier--Stokes produces such a cascade. It shows why the energy estimate alone
cannot rule one out.

In order for a singularity to form, velocity has to blow up. In order for
that to happen, its acceleration would also have to blow up, which means its
jerk has to blow up, and so on, and so on:
\[
u,\qquad \partial_tu,\qquad \partial_t^2u,\qquad \ldots
\]
What you'd expect to see in this pattern, if a singularity were to form, as
you go up the time derivatives, you'd expect to see less spatial regularity,
not more.

The equation moves in the other direction. Its first three time derivatives
are
\[
\begin{aligned}
\partial_tu
={}&\nu\Delta u
 -(u\!\cdot\nabla)u
 -\nabla p,\\[3pt]
\partial_t^2u
={}&\nu\Delta\partial_tu
 -(\partial_tu\!\cdot\nabla)u
 -(u\!\cdot\nabla)\partial_tu
 -\nabla\partial_tp,\\[3pt]
\partial_t^3u
={}&\nu\Delta\partial_t^2u
 -(\partial_t^2u\!\cdot\nabla)u
 -2(\partial_tu\!\cdot\nabla)\partial_tu
 -(u\!\cdot\nabla)\partial_t^2u
 -\nabla\partial_t^2p.
\end{aligned}
\]
\begin{samepage}
The pattern is
\[
\partial_t^{\,n+1}u
=\nu\Delta\partial_t^{\,n}u
-\sum_{a=0}^{n}\binom na
  \bigl(\partial_t^{\,a}u\!\cdot\nabla\bigr)
  \partial_t^{\,n-a}u
-\nabla\partial_t^{\,n}p.
\]
\end{samepage}
\begin{samepage}
Set
\[
V_n:=\partial_t^{\,n}u,
\qquad
Q_n:=\partial_t^{\,n}p,
\]
Taking the divergence and using incompressibility gives the pressure at that
order:
\[
-\Delta Q_n
=\sum_{a=0}^{n}\binom na
  \partial_i\partial_j
  \bigl((V_a)_i(V_{n-a})_j\bigr).
\]
\end{samepage}
The line to keep in view is \(V_{n+1}=\nu\Delta V_n-\cdots\). Moving one step
up in time brings in two spatial derivatives of the response below it.
Transport and pressure stay attached, but viscosity supplies the link from
temporal order to spatial Sobolev order.

At the critical height, set
\[
\Lambda=(-\Delta)^{1/2},
\qquad
E_s(t):=\frac12\|\Lambda^su(t)\|_2^2,
\qquad
H(t):=E_{1/2}(t).
\]
If \(\mathcal P_s\) denotes the complete transport--pressure contribution at
height \(s\), the energy identity is
\[
E_s'=\mathcal P_s-2\nu E_{s+1}.
\]
Starting from \(H=E_{1/2}\) and repeating that identity gives
\[
H^{(n)}
=(-2\nu)^nE_{n+1/2}
 +\sum_{j=0}^{n-1}(-2\nu)^j
   \partial_t^{\,n-1-j}\mathcal P_{j+1/2}.
\]
The first term has reached spatial height \(n+\tfrac12\); the sum carries every
transport--pressure term generated on the way. This is the opposite of the
blow-up picture. Rising through that temporal sequence was supposed to reveal
less spatial regularity. The equation exposes higher Sobolev height instead. As
you do that, you get higher order Sobolev embeddings:
\[
s>k+\frac32
\quad\Longrightarrow\quad
H^s(\mathbb R^3)\hookrightarrow C^k(\mathbb R^3),
\qquad
\bigcap_{s<\infty}H^s(\mathbb R^3)
=H^\infty(\mathbb R^3)
\hookrightarrow C^\infty(\mathbb R^3).
\]

Seeing a higher energy in the identity is not enough. It has to remain finite
all the way to \(T_*\). Fix smooth divergence-free finite-energy data \(u_0\),
viscosity \(\nu>0\), and its maximal pressure-complete solution \((u,p)\) on
\([0,T_*)\). Assume \(T_*<\infty\). For finite \(N\) and \(M\), collect the
adjacent rows in
\[
\mathfrak R_{N,M}(u;T)
:=\sum_{n=0}^{N}\left(
  \|V_n\|_{L^2(0,T;H^{M+1})}^2
 +\|V_{n+1}\|_{L^2(0,T;H^{M-1})}^2
\right).
\]
Adding the corresponding \(Q_n\) and \(\nabla Q_n\) coordinates gives the
finite pressure-complete rectangle \(\mathcal R_{N,M}[u_0;T]\). The
manuscript's whole-terminal theorem bounds each adjacent pair on the full
interval:
\[
V_n\in L^2(0,T_*;H^{M+1}),
\qquad
V_{n+1}\in L^2(0,T_*;H^{M-1}),
\qquad 0\le n\le N.
\]
Restriction to \((0,T)\) can only decrease those norms, and there are only
finitely many rows, so
\[
\sup_{0<T<T_*}\mathfrak R_{N,M}(u;T)<\infty.
\]
\begin{samepage}
The bound may depend on \(N\) and \(M\). No uniform constant over all orders is
needed. At \(n=0\), for every finite \(m\), the adjacent pair gives
\[
u\in L^2(0,T_*;H^{m+1}),
\qquad
\partial_tu\in L^2(0,T_*;H^{m-1}),
\]
The Lions--Magenes trace theorem applied to
\(H^{m+1}\hookrightarrow H^m\hookrightarrow H^{m-1}\), gives a unique
trace at the endpoint:
\[
u\in C([0,T_*];H^m)
\qquad\text{for every finite }m.
\]
\end{samepage}

\begin{samepage}
Increasing \(N\) or \(M\) does not create a new solution. Every rectangle is a
finite part of the maximal history from \(u_0\), so overlapping rectangles
agree on every shared coordinate. Their projective limit is
\[
\mathcal I[u_0]
:=\varprojlim_{N,M<\infty}\mathcal R_{N,M}[u_0;T_*],
\qquad
u\in\bigcap_{r,m<\infty}H^r(0,T_*;H^m(\mathbb R^3)).
\]
\end{samepage}
The intersection is not one bound over infinitely many orders. Each coordinate
was bounded at a finite \((N,M)\); compatibility says that all of those finite
bounds belong to one velocity--pressure history.

The zeroth temporal row, \(\mathscr S_0:=\pi_0\mathcal I[u_0]\), now has one
terminal value at every finite spatial height:
\[
u_*:=u(T_*)
\in\bigcap_{m<\infty}H^m_\sigma(\mathbb R^3)
=H^\infty_\sigma(\mathbb R^3)
\hookrightarrow C^\infty(\mathbb R^3).
\]
The recurrence and pressure equation recover every finite coordinate of the
terminal jet from \(u_*\). The uniform \(H^3\) row gives a local existence time
\(\delta_0>0\) that works for \(u_*\) and for every sufficiently late slice
\(u(t_0)\). Choose \(t_0<T_*\) with \(T_*-t_0<\delta_0/2\), and restart from
\(u(t_0)\). Uniqueness makes this restart identical to the original solution
on \([t_0,T_*)\), so it reaches \(T_*\) with value \(u_*\). To the right of
\(T_*\), uniqueness identifies it with the solution restarted from \(u_*\).
The two restarts join into one smooth pressure-complete solution beyond
\(T_*\), contradicting maximality. Thus \(T_*=\infty\).

Once \(V_r\in C_tH_x^m\) is available for every finite \(r\) and \(m\), any
finite downstream estimate comes from choosing
\[
m>k+\frac32-\frac3p,
\qquad
2\le p\le\infty,
\]
which gives, on every finite time interval,
\[
V_r\in C_tH_x^m
\hookrightarrow L_t^qW_x^{k,p},
\qquad 1\le q\le\infty.
\]

\[
u_0
\longrightarrow \{(V_n,Q_n)\}_{n\ge0}
\longrightarrow \{\mathcal R_{N,M}[u_0]\}_{N,M<\infty}
\longrightarrow \mathcal I[u_0]
\longrightarrow \mathscr S_0
\longrightarrow u_*\in H^\infty
\longrightarrow T_*=\infty.
\]
